% chapters/03_schnittstellen.tex

\chapter{Schnittstellen}
\label{chap:schnittstellen}
\label{chap:umsetzung}
\label{chap:methoden}

Die Tabellen fassen \doc{docs/schnittstellen.md} zusammen; dort steht zu jeder Angabe die
Quelldatei im Code.

\section{Datensatz}
\label{sec:if-datensatz}
\label{sec:datensatz}
\label{sec:training}

\begin{table}[htb]
  \centering
  \begin{tabularx}{\linewidth}{lX}
    \toprule
    Quelle & \texttt{unitreerobotics/G1\_Dex3\_BlockStacking\_Dataset} (Hugging Face), LeRobot
             v3.0~\cite{lerobot2024}, 301 Episoden, 281\,196 Frames, 30~Hz \\
    Konvertierung & beim Containerstart nach LeRobot v2.1, Ablage
             \texttt{/data/unitreerobotics/\dots}; \texttt{meta/modality.json} aus
             \texttt{examples/G1\_DEX3/modality\_4cam.json} \\
    Bilder & \texttt{cam\_left\_high}, \texttt{cam\_right\_high}, \texttt{cam\_left\_wrist},
             \texttt{cam\_right\_wrist}; RGB, 640$\times$480 \\
    Zustand, Aktion & je 28 Werte (rad): \texttt{left\_arm} 0--6, \texttt{right\_arm} 7--13,
             \texttt{left\_dex3} 14--20, \texttt{right\_dex3} 21--27 \\
    Aktionsform & Arme relativ zum Zustand, Hände absolut; Block von 16 Schritten \\
    Sprache & \texttt{annotation.human.task\_description}, Standardtext
             \enquote{stack three block} \\
    Split & \texttt{TRAIN\_TEST\_SPLIT=1}: Episoden 0--239 Training, 240--300 Test
             (Anteil 0,8 nach Episodenindex, ohne Zufall) \\
    \bottomrule
  \end{tabularx}
  \caption{Datensatz-Schnittstelle. Gelenkreihenfolge: \cref{anh:datensatz}.}
  \label{tab:if-datensatz}
\end{table}

\section{Checkpoint}
\label{sec:if-checkpoint}

\begin{table}[htb]
  \centering
  \begin{tabularx}{\linewidth}{lX}
    \toprule
    Pfad & \texttt{/data/g1\_dex3\_finetune/<lauf>/<experiment>/checkpoint-<N>}, z.\,B.\
           \texttt{blockstacking\_vision/g1\_dex3\_blockstacking\_vision\_v1}; N1.7 mit Suffix
           \texttt{\_n17} \\
    Takt & lokal alle 2000 Schritte (höchstens 15 behalten), KISSKI alle 5000 (40) \\
    Inhalt & \texttt{config.json}, Gewichte (safetensors), \texttt{experiment\_cfg/},
             \texttt{processor/}, \texttt{wandb\_config.json}, Optimizer-Zustand \\
    Version & \texttt{model\_type} in \texttt{config.json}: \texttt{Gr00tN1d6} oder
              \texttt{Gr00tN1d7}; \texttt{GROOT\_VERSION=auto} liest ihn aus \\
    Ablage & Hugging Face (Projekt-Repos öffentlich), ohne Optimizer:
             \texttt{upload\_checkpoint.py --checkpoint <dir> --repo
             luca-mue/groot-g1dex3-<lauf>-ckpt<N>} \\
    \bottomrule
  \end{tabularx}
  \caption{Checkpoint-Schnittstelle.}
  \label{tab:if-checkpoint}
\end{table}

\section{Policy-Server}
\label{sec:if-server}
\label{sec:simulation}

Der Server lädt einen Checkpoint und beantwortet Anfragen über \ac{zmq} (REQ/REP,
\doc{tcp://<host>:5555}). Nachrichten sind mit msgpack kodiert, numpy-Arrays als
\doc{{"__ndarray_class__": true, "as_npy": <Bytes von np.save>}}. Eine Anfrage hat die Form
\doc{{"endpoint": <name>, "data": {...}}}; Endpunkte sind \texttt{ping},
\texttt{reset}, \texttt{get\_action}, \texttt{get\_modality\_config} und \texttt{kill}. Fehler
kommen als \doc{{"error": <text>}} zurück.

\begin{table}[htb]
  \centering
  \begin{tabularx}{\linewidth}{lX}
    \toprule
    Anfrage \texttt{get\_action} & \texttt{video.<kamera>}: \texttt{(1,1,480,640,3)} uint8 für
      alle vier Kameras; \texttt{state.<gruppe>}: \texttt{(1,1,7)} float32 für die vier
      Gruppen; \texttt{annotation.human.task\_description}: Tupel mit einem String \\
    Antwort & \texttt{[action, info]}, \texttt{action.<gruppe>}: \texttt{(1,16,7)} float32,
      zusammen 16 absolute Gelenkziele $\times$ 28 \\
    Ausführung & 30~Hz; von 16 Schritten werden \texttt{EXECUTION\_HORIZON} = 8 ausgeführt,
      dann neu angefragt; Budget $8/30$~s = 267~ms \\
    Zeitlimits & Sim-Client 20~s je Anfrage; nach Timeout REQ-Socket neu aufbauen \\
    \bottomrule
  \end{tabularx}
  \caption{Protokoll des Policy-Servers. Ein lauffähiger Beispiel-Client steht in
    \doc{docs/schnittstellen.md}, Abschnitt~3.}
  \label{tab:if-server}
\end{table}

Der Server rechnet die relativen Arm-Aktionen mit dem mitgeschickten Zustand in absolute
Ziele um. Ein Client muss deshalb immer den aktuellen Zustand senden.

\section{Sim-Eval}
\label{sec:if-sim}
\label{sec:validierung}
\label{sec:eval-methodik}

\paragraph{Eingaben.} Eine Layout-Datei legt je Testepisode die drei Würfel fest
(\texttt{episodes.<nr>.cubes}: $x,y$ in m, rot/grün/gelb; \texttt{cubes\_yaw\_deg}; dazu
\texttt{n\_train}). Die Lagen werden aus dem ersten Bild der echten Episode bestimmt (Farbe,
Strahl auf die Tischebene; Fehler 0,3--0,4~cm gegen bekannte Lagen). Standard der Messungen ist
\texttt{layout\_heldout\_dy3.json}: Testlagen, Würfel je Seite 3~cm nach außen
(\cref{sec:erg-registrierung}). Die Szene aus Lauf~70 setzt zusätzlich Tischgröße und
Handgelenkkamera (\texttt{SIM\_TABLE\_SIZE}, \texttt{SIM\_WRIST\_*}).

\paragraph{Ausgaben.} Je Lauf eine JSON-Datei \texttt{results\_<zeit>\_<checkpoint>.json}
mit Kopf (Checkpoint, Layout, Parameter, 95-\%-Intervalle) und einem Eintrag je Episode.
Abgebrochene Episoden zählen in keiner Rate.

\begin{table}[htb]
  \centering
  \begin{tabularx}{\linewidth}{lX}
    \toprule
    Stufe & Bedingung \\
    \midrule
    \texttt{reached} & Fingerkuppe $\le$ 4~cm von der Würfelmitte \\
    \texttt{grasp\_attempted} & Fingerspanne des Kommandos $\ge$ 30\,\% der Demonstrationsspanne
      von 2,094~rad \\
    \texttt{touched} & Würfel um $>$ 1~cm verschoben \\
    \texttt{lifted} & Würfelschwerpunkt $>$ 2~cm angehoben (zählt auch Kippen) \\
    \texttt{lifted\_clear} & Würfelunterseite $>$ 2~cm über dem Tisch \emph{und} im selben
      Schritt eine Fingerkuppe $\le$ 4~cm entfernt \\
    \texttt{stacked} & drei Würfel horizontal $<$ 3~cm versetzt, Höhendifferenz der
      Würfelmitten $>$ 8~cm, alle $<$ 0,05~m/s \\
    \bottomrule
  \end{tabularx}
  \caption{Stufen der Closed-Loop-Bewertung. Jede Stufe wird für sich geprüft.
    \texttt{lifted\_clear} kam am 13.10.\ hinzu, weil \texttt{lifted} Kippen mitzählt;
    die übrigen Schwellen standen vor den Messungen fest.}
  \label{tab:if-stufen}
\end{table}

\paragraph{Visuelles Gate.} Ohne Isaac~Sim: Die Policy sagt teacher-forced Aktionen auf
gerenderten Prüfbildern (4 Episoden, je 600 Schritte, Würfel 5~cm und 4~cm) und auf den
echten Bildern derselben Episoden vorher. Kennzahl ist der Median der \ac{mse} über die
Hand-Dimensionen, gemittelt über die Seeds 0, 1 und 2 (Streuung etwa $\pm$5\,\%).
Bestanden ist das Gate, wenn die Kennzahl auf beiden Prüfsets unter 80\,\% des Werts des
Ausgangsmodells derselben Generation ohne Co-Training liegt (N1.6 vision-v2 bzw.\ N1.7
vision-v1, je @44000, Lauf~80) und auf Echtbildern nicht deutlich steigt. Ergebnis:
\texttt{visgate\_*.json}.

\section{Container und Einstiegspunkte}
\label{sec:if-container}
\label{sec:docker}

\begin{table}[htb]
  \centering
  \begin{tabularx}{\linewidth}{lX}
    \toprule
    Training & Image \texttt{lucam03/projekt-humanoider-roboter:latest} (N1.6) bzw.\
      \texttt{:latest-n17}; Start lokal \texttt{Training/setup\_and\_train\_dockerhub\_pull.sh},
      KISSKI \texttt{sbatch --export=ALL Training/kisski\_submit.sh} \\
    Simulation & Image \texttt{lucam03/projekt-humanoider-roboter-sim-standalone:latest}
      (Isaac~Sim~6.0); alle Aktionen über \texttt{Simulation/server\_rl\_run.sh <aktion>} \\
    Menü & \texttt{./run.sh} führt zu allen Aktionen und zeigt den direkten Befehl \\
    Daten & Training \texttt{/data} im Container (KISSKI: Projektspeicher); Sim
      \texttt{\textasciitilde/groot-rl-data} = \texttt{/data} im Container \texttt{groot-rl} \\
    Variablen & \texttt{HF\_TOKEN} (Pflicht), \texttt{GROOT\_VERSION}, \texttt{MAX\_STEPS},
      \texttt{GLOBAL\_BATCH\_SIZE}, \texttt{NUM\_GPUS}, \texttt{TUNE\_VISUAL},
      \texttt{USE\_COTRAIN}, \texttt{TRAIN\_TEST\_SPLIT}; vollständig in
      \doc{docs/training/env-vars.md} \\
    \bottomrule
  \end{tabularx}
  \caption{Container, Einstiegspunkte und Konfiguration.}
  \label{tab:if-container}
\end{table}

Der Trainings-Container hält alle Daten; er darf gestoppt, aber nicht gelöscht werden
(kein \texttt{docker run --rm}).

\section{Echter Roboter}
\label{sec:if-roboter}

Eine Anbindung an den G1 fehlt. Eine Roboter-Seite müsste die vier Kamerabilder und die 28
Gelenkwinkel lesen, sie nach \cref{tab:if-server} an den Server schicken und die
zurückkommenden Gelenkziele mit 30~Hz an Arme und Hände geben. Gleichgewicht und Beine
regelt die Policy nicht.
